\part*{Prerrequisitos}

Estos apuntes están escritos de forma que los prerrequisitos para entenderlos son sorprendentemente laxos para el nivel de material que se cubre. Hay que saber lo siguiente:
\begin{itemize}
  \item Programación básica en algún lenguaje de programación, incluyendo recursión. 
  \begin{itemize}
    \item El seudocódigo que se emplea es bastante parecido a Python en la sintaxis; sin embargo, es legible si se tiene experiencia con cualquier lenguaje de programación imperativo.
  \end{itemize}
  \item Algo de teoría sobre las estructuras de datos más básicas: arreglos, listas, etc. En qué difieren y cuáles son las ventajas de una sobre otra.
  \item Matemática básica
  \begin{itemize}
    \item Notación matemática básica: conjuntos, operadores, cuantificadores, etc.
    \item Aritmética básica: sumar, multiplicar, potenciar, etc.
    \item Álgebra elemental: despejar variables, factorizar, etc.
    \item Funciones elementales: logaritmos, exponenciales, polinomios, etc.
  \end{itemize}
\end{itemize}
% TODO: Verificar